% Irodov Part III §3.2, original 54--65, PDF pp.108--109,
% printed pp.99--100. Image-charge signs checked; conductor assumed ideal.
\electroconductproblem{54}{99}
一个小球由劲度系数为 $k_s$ 的绝缘弹性细线
悬挂在无限大水平导体平面上方。
小球充电后比原平衡位置向下降
$x$，此时距导体平面 $\ell$。
求小球电荷模。
\exsolution
无限导体平面用镜像法，
在平面下方距平面同为 $\ell$ 的位置
放置大小相同、异号的镜像电荷。
两电荷相距 $2\ell$，
作用于真小球的镜像吸引力为
\[
F_e=\frac{q^2}{4\pi\varepsilon_0(2\ell)^2}
=\frac{q^2}{16\pi\varepsilon_0\ell^2}.
\]
充电前后球的重力相同；
新增加的弹性线拉力为 $k_sx$，
故 $F_e=k_sx$。
因此
\[
\boxed{|q|=4\ell\sqrt{\pi\varepsilon_0k_sx}}.
\]
\excheck $x$ 若以 cm 给出，
代数值时须先换算成 m；
结果与平面是否带净电荷无关的前提是
理想接地或足够大的导体平面。

\electroconductproblem{55}{99}
点电荷 $q$ 位于距无限大接地导体平面
$\ell$ 处。
求把该电荷缓慢移到无穷远需要做的外功。
\exsolution
镜像电荷为 $-q$，位于平面另一侧
对称位置。真电荷距平面 $z$ 时，
受到指向平面的吸引力
$F_z=-kq^2/(2z)^2=-q^2/(16\pi\varepsilon_0z^2)$。
为准静态地从 $z=\ell$ 移至无穷远，
外力必须做
\[
W_{\rm ext}=-\int_\ell^\infty F_z\,\mathrm dz
=\frac{q^2}{16\pi\varepsilon_0}
\int_\ell^\infty z^{-2}\,\mathrm dz.
\]
故
\[
\boxed{W_{\rm ext}=\frac{q^2}{16\pi\varepsilon_0\ell}}.
\]
\excheck 不应把真电荷与镜像电荷距离当作固定；
且镜像电荷随真电荷移动，
直接用“两真实点电荷的势能”会有
多出一倍的计数问题。

\electroconductproblem{56}{99}
两个真点电荷 $+q,-q$ 在同一高度
$h=\ell/2$、相互水平距离为 $\ell$，
二者均位于无限接地导体平面上方。
求：(a) 每个真电荷受力模；
(b) 两电荷中点处电场强度模。
\exsolution
取平面为 $z=0$，真电荷位置为
$(0,h),(\ell,h)$。
对应镜像分别为 $-q$ 于 $(0,-h)$、
$+q$ 于 $(\ell,-h)$。
(a) 真正的 $+q$ 受三项作用：
真 $-q$ 的水平吸引
$+kq^2/\ell^2$，
自身镜像的向下吸引
$-kq^2/\ell^2$，
另一镜像 $+q$ 的斜向排斥
$kq^2/(2\ell^2)$、方向 $(-1,1)/\sqrt2$。
合力分量互为相反数：
\[
F_x=\frac{kq^2}{\ell^2}
\left(1-\frac{1}{2\sqrt2}\right),\quad F_z=-F_x.
\]
所以
\[
\boxed{F=\frac{(2\sqrt2-1)q^2}
{8\pi\varepsilon_0\ell^2}}.
\]
(b) 中点两真电荷电场沿连线同向叠加，
贡献 $8kq/\ell^2$；
两个镜像对该方向的合贡献为
$-8kq/(5\sqrt5\,\ell^2)$，
竖直分量相消。故
\[
\boxed{E=
\frac{2|q|}{\pi\varepsilon_0\ell^2}
\left(1-\frac1{5\sqrt5}\right)}.
\]
\excheck 所有镜像仅为解边界值的数学工具，
不是真实可自由移动的第三、第四个点电荷。

\electroconductproblem{57}{99}
两个无限大、相互垂直的接地导体半平面形成直角，
一点电荷 $q$ 位于两平面之间，
距每个平面均为 $\ell$。
求该电荷所受合电场力模。
\exsolution
在直角区域坐标中取真电荷 $(\ell,\ell)$。
两个单次反射镜像电荷为 $-q$，
位置分别在 $(-\ell,\ell)$、
$(\ell,-\ell)$；
双次反射镜像为 $+q$，
位置在 $(-\ell,-\ell)$。
前两个分别沿 $-x,-y$ 方向吸引，
每个力模为 $kq^2/(4\ell^2)$。
第三个沿 $(+x,+y)$ 排斥，
力模为 $kq^2/(8\ell^2)$，
每分量为该值除以 $\sqrt2$。
所以
\[
F_x=F_y=
-\frac{kq^2}{4\ell^2}
+\frac{kq^2}{8\sqrt2\,\ell^2}.
\]
合力模
\[
\boxed{F=
\sqrt2|F_x|=
\frac{(2\sqrt2-1)q^2}
{32\pi\varepsilon_0\ell^2}}.
\]
\excheck 必须有\emph{三个}镜像，
且双次反射镜像与真电荷同号；
漏掉它会造成错误的系数。

\electroconductproblem{58}{99}
点电偶极矩为 $\bm p$ 的偶极子
距接地无限大导体平面 $\ell$，
偶极矩垂直于导体表面。
求作用在偶极子上的力模。
\exsolution
垂直导体平面的偶极子的镜像偶极矩
与原偶极矩\emph{同向}，
两偶极子中心距为 $2\ell$。
在两同轴同向偶极子的轴线上，
相互作用势能函数的距离部分为
$U=-2kp^2/d^3$。
对\emph{真偶极子在镜像场中的力}，
在固定镜像构造下求其所受局部场梯度：
\[
F=\frac{6kp^2}{d^4}\bigg|_{d=2\ell}
=\frac{6kp^2}{16\ell^4}.
\]
因此
\[
\boxed{F=\frac{3p^2}{32\pi\varepsilon_0\ell^4}}.
\]
方向朝向导体平面。
\excheck 能量法处理镜像随距离改变时，
镜像系统能量要乘 $1/2$；
直接用 $\bm F=(\bm p\cdot\nabla)\bm E_{\rm image}$
不会混淆这一因子。

\electroconductproblem{59}{99}
点电荷 $q$ 位于无限接地导体平面上方
高度 $\ell$。
求导体表面，距电荷垂足水平距离为 $r$
处的感应面电荷密度，并求总感应电荷。
\exsolution
取导体为 $z=0$，镜像 $-q$ 在 $z=-\ell$。
在平面上方，
两者的法向电场分量叠加为
\[
E_z(0^+,r)=
-\frac{q\ell}{2\pi\varepsilon_0
(r^2+\ell^2)^{3/2}}.
\]
导体内部 $\bm E=0$，
表面法向边界条件
$\sigma=\varepsilon_0E_z(0^+,r)$，
于是
\[
\boxed{\sigma(r)=
-\frac{q\ell}{2\pi(r^2+\ell^2)^{3/2}}}.
\]
积分整个平面：
\[
Q_{\rm ind}
=\int_0^\infty\sigma(r)2\pi r\,\mathrm dr
=-q\ell\int_0^\infty
\frac{r\,\mathrm dr}{(r^2+\ell^2)^{3/2}}
=\boxed{-q}.
\]
\excheck 对正点电荷，感应面电荷密度应全为负；
式子积分总量为 $-q$ 是重要自校。

\electroconductproblem{60}{99}
一根无限长线电荷密度 $\lambda$ 的细线，
平行于无限接地导体平面，距平面 $\ell$。
求 (a) 每单位长度所受电场力模；
(b) 平面上距细线所在垂直投影平面
横向距离 $x$ 的感应面电荷密度。
\exsolution
用镜像细线，线密度为 $-\lambda$，
平行放在平面另一侧，距真线 $2\ell$。
两无限直线之间每单位长度的静电作用力
\[
\boxed{\frac FL=
\frac{\lambda^2}{2\pi\varepsilon_0(2\ell)}
=\frac{\lambda^2}{4\pi\varepsilon_0\ell}}.
\]
方向朝向导体平面。
(b) 设导体面为 $z=0$，
真线位于 $(x,z)=(0,\ell)$，
镜像线位于 $(0,-\ell)$。
面上法向场
\[
E_z(0^+,x)
=-\frac{\lambda\ell}{\pi\varepsilon_0(x^2+\ell^2)}.
\]
故
\[
\boxed{\sigma(x)=
-\frac{\lambda\ell}{\pi(x^2+\ell^2)}}.
\]
\excheck 对正 $\lambda$，感应面电荷为负，
且沿每单位线长积分
$\int_{-\infty}^\infty\sigma(x)\,\mathrm dx=-\lambda$。

\electroconductproblem{61}{100}
半无限长直细线垂直于无限接地导体平面，
其近端距平面为 $\ell$，
线电荷密度为 $\lambda$。
设该线在平面上的垂足为 O，
求 (a) O 点感应电荷面密度；
(b) 与 O 距离 $r$ 处的分布。
\exsolution
设真线沿 $z\ge\ell$ 延伸，
每个高度 $z$ 的电荷元
$\mathrm dq=\lambda\,\mathrm dz$，
其镜像在 $-z$。
由第59题的点电荷镜像诱导密度叠加：
\[
\sigma(r)=
-\frac{\lambda}{2\pi}
\int_\ell^\infty
\frac{z\,\mathrm dz}{(r^2+z^2)^{3/2}}.
\]
积分原函数为
$-1/\sqrt{r^2+z^2}$，
因此
\[
\boxed{\sigma(r)=
-\frac{\lambda}{2\pi\sqrt{r^2+\ell^2}}}.
\]
在 $r=0$：
\[
\boxed{\sigma(0)=-\frac{\lambda}{2\pi\ell}}.
\]
\excheck 本题的半无限线总电荷无限，
因而不要求整个平面的感应电荷积分有限；
面密度随 $r$ 缓慢衰减。

\electroconductproblem{62}{100}
半径 $R$ 的均匀带电细圆环总电荷 $q$，
其所在平面平行于无限接地导体平面，
且距离为 $\ell$。
求 (a) 轴线投影点处的感应面电荷密度；
(b) 环心处的电场强度模和电势。
\exsolution
真实环在 $z=\ell$，
镜像为带 $-q$ 的相同圆环在 $z=-\ell$。
(a) 对底部投影点，
两圆环在导体平面上的法向场叠加，
\[
E_z(0)=
-\frac{q\ell}{2\pi\varepsilon_0(R^2+\ell^2)^{3/2}},
\]
所以
\[
\boxed{\sigma(0)=
-\frac{q\ell}{2\pi(R^2+\ell^2)^{3/2}}}.
\]
(b) 环心处真实圆环的自场为0，
镜像环位于轴向距离 $2\ell$ 处，
其场模为
\[
\boxed{E=
\frac{|q|(2\ell)}
{4\pi\varepsilon_0(R^2+4\ell^2)^{3/2}}}.
\]
电势由真实环及镜像环相加：
\[
\boxed{\varphi_{\rm center}=
\frac q{4\pi\varepsilon_0}
\left(\frac1R-\frac1{\sqrt{R^2+4\ell^2}}\right)}.
\]
\excheck 电场方向指向导体面（当 $q>0$）；
电势是标量总和，必须计入真实环自电势。

\electroconductproblem{63}{100}
一颗总电荷为零、半径为 $R$ 的
\emph{孤立}导体球，
球心外距 $\ell>R$ 处有点电荷 $q$。
求导体球的平衡电势。
\exsolution
接地导体球的镜像解是：
在距球心 $R^2/\ell$ 处的镜像电荷
$q'=-qR/\ell$，
它与外部点电荷共同保证球面电势为0。
现在球要求总电荷为0而非接地，
需在球心再加补偿镜像
$q''=-q'=qR/\ell$。
由于球心点电荷使球面各处电势
增加相同常数
$kq''/R$，
外部点电荷与偏心镜像在球面仍恰好抵消，
故
\[
\boxed{\varphi_{\rm sphere}
=\frac{q''}{4\pi\varepsilon_0R}
=\frac{q}{4\pi\varepsilon_0\ell}}.
\]
\excheck 球电势不依赖半径 $R$，
但诱导面电荷\emph{分布}依赖 $R$；
本题的球是孤立无净电荷，
不能直接当作接地导体球。

\electroconductproblem{64}{100}
一颗半径为 $R_1,R_2$（$R_1<R_2$）
的中空导体球壳原本不带净电荷。
点电荷 $q$ 位于球腔内部，
距球心 O 为 $r<R_1$。
求球心 O 处电势。
\exsolution
真实点电荷在 O 处产生电势 $kq/r$。
为使金属内部场为零，
内球面出现总感应电荷 $-q$，
外球面出现总感应电荷 $+q$。
尽管内球面的电荷分布可能不均匀，
全部位于距 O 恰为 $R_1$ 的球面上，
故其在 O 处总电势为
$k(-q)/R_1$。
外球面均匀总电荷 $+q$ 在 O 的电势为
$kq/R_2$。
因此
\[
\boxed{\varphi(O)=
\frac q{4\pi\varepsilon_0}
\left(\frac1r-\frac1{R_1}+\frac1{R_2}\right)}.
\]
\excheck 内表面\emph{不是}均匀感应电荷，
但对球心电势，每个表面元距离相同，
故只需其总量。

\electroconductproblem{65}{100}
两个同心薄导体球，
内球半径为 $a$、带电 $q_1>0$，
外球半径为 $b>a$。
应给外球加多少总电荷 $q_2$，
才能让内球电势为零？
求各区间电势 $\varphi(r)$ 并描述图线。
\exsolution
内球表面电势为
$\varphi(a)=k(q_1/a+q_2/b)$，
故令其为0得到
\[
\boxed{q_2=-\frac ba q_1}.
\]
静电平衡时两球外电势按包围电荷求：
$r\le a$ 内球内部为常势零；
$a\le r\le b$，
由内球场积分并使用 $\varphi(a)=0$：
$\varphi(r)=kq_1(1/r-1/a)$；
$r\ge b$，
两球等效总电荷
$q_1+q_2=q_1(1-b/a)$，
所以
\[
\boxed{\varphi(r)=\frac{q_1}{4\pi\varepsilon_0}
\begin{cases}
0,&0\le r\le a,\\
\dfrac1r-\dfrac1a,&a<r<b,\\
\dfrac{1-b/a}{r},&r\ge b.
\end{cases}}
\]
电势图像在 $r=a$ 从0起向负值下降，
到 $r=b$ 后以负的 $1/r$ 形式趋向0，
且在两个金属表面电势连续。
\excheck 外球带电绝对值
$|q_2|>q_1$，故远处总电荷为负；
题目并未要求整个系统中性。
